% META: {"id":"1SPE-DERGLOBAL-EX-016","chapitre":"1SPE-DERIVATION-GLOBAL","type_objet":"exercice","capacites":["1SPE-DERIVATION-GLOBAL-C2"],"capacites_codes":["C2"],"methodes":["M2"],"parcours":2,"competences":["calculer","modeliser"],"duree_min":12,"mode_creation":"ex_nihilo","sources_inspiration":[],"fichier_tex":"chapitres/1SPE-DERIVATION-GLOBAL/exercices/1SPE-DERGLOBAL-EX-016.tex","corrige_tex":"chapitres/1SPE-DERIVATION-GLOBAL/corriges/1SPE-DERGLOBAL-CO-016.tex","status":"generated"}
% BEGIN-VERIFY
% from sympy import symbols, diff, simplify, solve, sqrt, Rational
% x = symbols('x')
% f = (x**2 - 2*x)/(x**2 + 1)
% fp = diff(f, x)
% num = 2*(x**2 + x - 1)
% assert simplify(fp - num/(x**2+1)**2) == 0
% sols = solve(fp, x)
% assert len(sols) == 2
% assert set(sols) == {-Rational(1,2) + sqrt(5)/2, -sqrt(5)/2 - Rational(1,2)}
% END-VERIFY
\begin{exercice}{1SPE-DERGLOBAL-EX-016}{2}{12}
Soit $f(x) = \dfrac{x^2-2x}{x^2+1}$ définie sur $\mathbb{R}$.
\begin{enumerate}
  \item Calculer $f'(x)$ et l'écrire sous la forme $\dfrac{2(x^2 + x - 1)}{(x^2+1)^2}$.
  \item Résoudre $f'(x) = 0$ (on acceptera une forme avec des radicaux).
  \item En déduire le tableau de variations de $f$.
\end{enumerate}
\end{exercice}
